\documentclass[11pt,a4paper]{article}
\usepackage[utf8]{inputenc}
\usepackage[vietnamese]{babel}
\usepackage{amsmath,amssymb}
\usepackage[version=4]{mhchem}
\usepackage{geometry}
\geometry{top=1.5cm,bottom=1.2cm,left=1.4cm,right=1.4cm,headheight=28pt,headsep=12pt,footskip=18pt}
\usepackage{graphicx}
\usepackage{xcolor}
\usepackage{tabularx}
\usepackage{array}
\usepackage{colortbl}
\usepackage{tikz}
\usepackage{fancyhdr}
\usepackage{enumitem}
\usepackage{multido}
\usepackage{tcolorbox}

\definecolor{hfRed}{HTML}{D32F2F}
\definecolor{hfGreen}{HTML}{8BC390}
\definecolor{mainGreen}{HTML}{2E7D32}
\definecolor{solBlue}{HTML}{1565C0}
\definecolor{tableGreen}{HTML}{E8F5E9}

\usepackage{ex_test}

\pagestyle{fancy}
\fancyhf{}
\renewcommand{\headrulewidth}{0.75pt}
\renewcommand{\headrule}{\hbox to\headwidth{\color{hfGreen}\leaders\hrule height \headrulewidth\hfill}}
\renewcommand{\footrulewidth}{0.75pt}
\renewcommand{\footrule}{\hbox to\headwidth{\color{hfGreen}\leaders\hrule height \footrulewidth\hfill}}

\fancyhead[L]{\small\itshape\color{hfRed}Tài liệu lưu hành nội bộ}
\fancyhead[R]{\small\color{hfRed}Thầy \textbf{TRẦN VĂN THẠNH} - 0777.470.803}

\fancyfoot[L]{\footnotesize\itshape\color{hfRed}CS1: 6/15 Nguyễn Hoàng, P. Kim Long, TP Huế \quad|\quad CS2: 24 Đặng Thái Thân, TP Huế}
\fancyfoot[R]{\small\textbf{\color{hfRed}Trang \thepage}}

\setlength{\parindent}{0pt}
\setlength{\parskip}{2pt}

\begin{document}

\begin{minipage}[t]{0.48\textwidth}
\textbf{SỞ GD - ĐT THÀNH PHỐ HUẾ}\\
\textbf{TRƯỜNG THPT ĐẶNG HUY TRỨ}\\
\textit{ĐỀ CHÍNH THỨC}\\
\textit{(Đề gồm có 03 trang)}
\end{minipage}
\hfill
\begin{minipage}[t]{0.48\textwidth}
\raggedleft
\textbf{\color{hfRed}KIỂM TRA CUỐI KỲ I}\\
\textbf{NĂM HỌC: 2025 -- 2026}\\
\textbf{Môn: HÓA HỌC -- Lớp 12}\\
\textit{Thời gian làm bài: 45 phút (không kể thời gian phát đề)}\\
\textbf{MÃ ĐỀ: 104}
\end{minipage}

\vspace{4pt}
\noindent
\textit{Họ và tên học sinh:} \dotfill \textit{Lớp:} \dotfill \textit{SBD:} \dotfill

\vspace{3pt}
\noindent
\textit{\footnotesize Cho NTK các nguyên tố: $\ce{Cu}=64;\ \ce{Ag}=108;\ \ce{C}=12;\ \ce{H}=1;\ \ce{O}=16;\ \ce{N}=14;\ \ce{K}=39;\ \ce{Na}=23;\ \ce{Zn}=65;\ \ce{Ni}=59;\ \ce{Sn}=119;\ \ce{Mg}=24$.}

\vspace{3pt}
\hrule color{hfGreen} height 0.5pt

\vspace{4pt}
\noindent
{\color{mainGreen}\textbf{\Large A. TRẮC NGHIỆM (7,0 điểm).}}

\vspace{3pt}
\noindent
{\color{mainGreen}\textbf{\large Phần I: Trắc nghiệm nhiều lựa chọn (4,0 điểm).}}\\
\textit{\small Học sinh trả lời từ câu 1 đến câu 16. Mỗi câu hỏi học sinh chỉ chọn một phương án.}
\vspace{2pt}

\begin{ex}
Cho các thế điện cực chuẩn của các cặp oxi hóa -- khử của kim loại như sau:
(1) $\ce{Ag+}/\ce{Ag}\ (= 0,80\text{V})$; (2) $\ce{Cu^2+}/\ce{Cu}\ (= 0,34\text{V})$; (3) $\ce{Na+}/\ce{Na}\ (-2,71\text{V})$; (4) $\ce{Mg^2+}/\ce{Mg}\ (-2,37\text{V})$.
Ion kim loại có tính oxi hoá mạnh nhất là ion kim loại trong cặp oxi hóa khử
\choice
{(2)}
{(4)}
{(3)}
{\True (1)}
\end{ex}
\begin{ex}
Chất nào dưới đây không phải là amino acid?
\choice
{Glycine}
{Lysine}
{Glutamic acid}
{\True Aniline}
\end{ex}
\begin{ex}
Chất nào sau đây có thể tham gia phản ứng trùng ngưng?
\choice
{$\ce{CH2=CHC6H5}$}
{\True $\ce{H2N[CH2]5COOH}$}
{$\ce{CH2=CHCOOH}$}
{$\ce{CH2=CH(CH3)}$}
\end{ex}
\begin{ex}
Một pin điện hoá có điện cực là Zn nhúng trong dung dịch $\ce{ZnSO4}$ và điện cực là Cu nhúng trong dung dịch $\ce{CuSO4}$. Sau một thời gian pin đó phóng điện thì
\choice
{khối lượng cả hai điện cực Zn và Cu đều giảm}
{khối lượng cả hai điện cực Zn và Cu đều tăng}
{khối lượng điện cực Zn tăng còn khối lượng điện cực Cu giảm}
{\True khối lượng điện cực Zn giảm còn khối lượng điện cực Cu tăng}
\end{ex}
\begin{ex}
Chất nào dưới đây không tan trong nước nhưng tan được trong dung dịch Schweizer?
\choice
{Saccharose}
{\True Cellulose}
{Maltose}
{Fructose}
\end{ex}
\begin{ex}
Cho tetrapeptide X: Gly-Ala-Glu-Val. Amino acid đầu N là
\choice
{Val}
{\True Gly}
{Glu}
{Ala}
\end{ex}
\begin{ex}
Hình dưới đây là ký hiệu của 6 polymer nhiệt dẻo phổ biến có thể tái chế:
\begin{center}
\includegraphics[width=9.5cm]{C:/Antigravity_Thanh/San_Pham/images_de104/ky_hieu_tai_che_6.png}
\end{center}
Các ký hiệu này thường được in trên bao bì, vỏ hộp, đồ dùng,... để giúp nhận biết vật liệu polymer cũng như thuận lợi cho việc thu gom, tái chế. Polymer có ký hiệu số 5 được điều chế bằng phản ứng trùng hợp monomer nào dưới đây?
\choice
{$\ce{CH2=CH2}$}
{$\ce{CH2=CH-Cl}$}
{\True $\ce{CH2=CH-CH3}$}
{$\ce{CH2=CH-C6H5}$}
\end{ex}
\begin{ex}
Cho các hợp chất có công thức cấu tạo dưới đây:
\begin{center}
\includegraphics[width=8.5cm]{C:/Antigravity_Thanh/San_Pham/images_de104/ctct_amino_acid_cau8.png}
\end{center}
Những hợp chất nào trong số các chất trên thuộc loại $\alpha$ -- amino acid?
\choice
{\True Chất (1), chất (3)}
{Chất (1) và chất (2)}
{Chất (1), chất (2) và chất (3)}
{Chất (2) và chất (3)}
\end{ex}
\begin{ex}
Khẳng định nào sau đây đúng?
\choice
{Chất béo là triester của glycerol với acid vô cơ}
{Chất béo là chất rắn tan tốt trong nước}
{\True Chất béo không tan trong nước nhưng tan nhiều trong dung môi hữu cơ}
{Dầu ăn và mỡ bôi trơn có cùng thành phần nguyên tố}
\end{ex}
\begin{ex}
Xà phòng có thành phần chính là
\choice
{muối sodium hoặc potassium của acid bất kì}
{glycerol và sodium chloride hoặc potassium chloride}
{\True muối sodium hoặc potassium của acid béo}
{muối sodium hoặc potassium của carboxylic acid}
\end{ex}
\begin{ex}
Nhận định nào sau đây đúng khi nói về amino acid?
\choice
{dễ nóng chảy và ít tan trong nước}
{\True có nhiệt độ nóng chảy cao và tan tốt trong nước}
{có nhiệt độ nóng chảy cao và ít tan trong nước}
{dễ nóng chảy và tan tốt trong nước}
\end{ex}
\begin{ex}
Phát biểu nào sau đây đúng?
\choice
{\True PVC được điều chế bằng phản ứng trùng hợp}
{Tơ visco thuộc loại tơ tổng hợp}
{Amylose có cấu trúc mạch phân nhánh}
{Tơ tằm thuộc loại tơ nhân tạo}
\end{ex}
\begin{ex}
Một tripeptide X được cấu thành từ 2 phân tử Ala và 1 phân tử Gly. Công thức cấu tạo của X không thể là
\choice
{\True Gly-Ala-Gly}
{Gly-Ala-Ala}
{Ala-Gly-Ala}
{Ala-Ala-Gly}
\end{ex}
\begin{ex}
Có thể phân biệt glucose và fructose bằng cách cho từng mẫu chất tác dụng với
\choice
{thuốc thử Tollens}
{\True nước bromine}
{dung dịch $\ce{NH3}$}
{$\ce{Cu(OH)2}$ trong môi trường kiềm, đun nóng}
\end{ex}
\begin{ex}
Cao su buna -- N (hay còn gọi là cao su nitrile, có kí hiệu là NBR) là loại cao su tổng hợp có khả năng chịu dầu mỡ tốt nên dùng làm ống dẫn nhiên liệu, gioăng làm kín trong các máy móc. Thực hiện phản ứng trùng hợp các chất nào dưới đây thu được sản phẩm cao su buna -- N?
\choice
{$\ce{CH2=CHCH=CH2}$ và $\ce{C6H5CH=CH2}$}
{$\ce{CH2=C(CH3)CH=CH2}$ và $\ce{CH2=CHCN}$}
{\True $\ce{CH2=CHCH=CH2}$ và $\ce{CH2=CHCN}$}
{$\ce{CH2=CHCH=CH2}$ và $\ce{N2}$}
\end{ex}
\begin{ex}
Tên gọi của $\ce{HCOOC2H5}$ là
\choice
{methyl acetate}
{ethyl acetate}
{\True ethyl formate}
{methyl formate}
\end{ex}
\vspace{4pt}
\noindent
{\color{mainGreen}\textbf{\large Phần II: Trắc nghiệm dạng Đúng/Sai (2,0 điểm).}}\\
\textit{\small Học sinh trả lời từ câu 1 đến câu 2. Trong mỗi ý a), b), c), d) ở mỗi câu, học sinh chọn đúng hoặc sai.}
\vspace{2pt}

\begin{ex}
Poly(ethylene terephthalate) (viết tắt là PET) là một polymer được điều chế từ terephthalic acid và ethylene glycol theo phương trình:
\begin{center}
\includegraphics[width=12.5cm]{C:/Antigravity_Thanh/San_Pham/images_de104/phan_ung_pet.png}
\end{center}
PET được sử dụng để sản xuất tơ, chai đựng nước uống, hộp đựng thực phẩm. Để thuận lợi cho việc nhận biết, sử dụng và tái chế thì các đồ nhựa làm từ vật liệu chứa PET thường được in kí hiệu như hình trên.
\choiceTF
{\True Trong một mắt xích PET, phần trăm khối lượng carbon là $62,5\%$}
{Phản ứng tổng hợp PET từ terephthalic acid và ethylene glycol thuộc loại phản ứng trùng hợp}
{\True Sản phẩm làm bằng nhựa số 01(PET) có thể được tái chế nhưng không nên tái sử dụng để đựng thực phẩm, nước uống}
{\True Tơ được chế tạo từ PET thuộc loại tơ tổng hợp}
\end{ex}
\begin{ex}
Trong phòng thí nghiệm, một học sinh nhúng thanh đồng có khối lượng $12,340\text{ g}$ vào $255\text{ mL}$ dung dịch $\ce{AgNO3}\ 0,125\text{M}$. Bằng quan sát, học sinh đó đã khẳng định có phản ứng xảy ra.
\choiceTF
{Tính oxi hóa của $\ce{Cu^2+}$ mạnh hơn $\ce{Ag+}$}
{Dung dịch từ không màu chuyển sang màu nâu đỏ}
{\True Phương trình hóa học của phản ứng: $\ce{Cu + 2AgNO3 -> Cu(NO3)2 + 2Ag}$}
{Giả thuyết toàn bộ lượng Ag giải phóng đều bám vào thanh đồng thì khối lượng thanh đồng sau khi phản ứng kết thúc là $13,35\text{ gam}$}
\end{ex}
\vspace{4pt}
\noindent
{\color{mainGreen}\textbf{\large Phần III: Trắc nghiệm trả lời ngắn (1,0 điểm).}}\\
\textit{\small Học sinh trả lời từ câu 1 đến câu 4.}
\vspace{2pt}

\begin{ex}
Cho các phát biểu:
\begin{enumerate}[label=(\alph*),leftmargin=20pt,itemsep=1pt]
\item Oxi hóa glucose bằng thuốc thử Tollens tạo ra gluconic acid.
\item Dung dịch của glucose hay saccharose đều có thể hòa tan $\ce{Cu(OH)2}$ ở điều kiện thường.
\item Cellulose trinitrate là nguyên liệu để sản xuất thuốc súng không khói.
\item Trong cellulose, các gốc $\beta$-glucose chỉ liên kết với nhau bằng liên kết $\beta$-1,4-glycoside.
\item Glucose và fructose có thể chuyển hóa lẫn nhau trong môi trường kiềm.
\item Glucose là hợp chất thuộc loại monosaccharide.
\end{enumerate}
Số phát biểu đúng là bao nhiêu?
\hfill \textbf{Trả lời:} \dotfill
\end{ex}
\begin{ex}
Cho các loại tơ: bông, tơ capron, tơ cellulose acetate, tơ tằm, len, tơ nitron, nylon-6,6. Có bao nhiêu tơ thuộc loại tơ tổng hợp?
\hfill \textbf{Trả lời:} \dotfill
\end{ex}
\begin{ex}
Xét pin Galvani hoạt động với phương trình tương ứng: $\ce{Zn(s) + Cu^2+(aq) -> Cu(s) + Zn^2+(aq)}$\\
Cho các phát biểu:
\begin{enumerate}[label=(\alph*),leftmargin=20pt,itemsep=1pt]
\item Điện cực Copper giảm khối lượng và điện cực Copper là cực dương.
\item Sức điện động của pin $E^\circ_\text{pin} = E^\circ(+) - E^\circ(-)$.
\item Điện cực Zinc giảm khối lượng và điện cực Zinc là cực âm.
\item Khi pin hoạt động dòng điện có chiều từ Copper sang Zinc.
\item Thế điện cực chuẩn của cặp $\ce{Zn^2+}/\ce{Zn}$ nhỏ hơn của cặp $\ce{Cu^2+}/\ce{Cu}$.
\end{enumerate}
Có bao nhiêu phát biểu đúng trong các phát biểu trên?
\hfill \textbf{Trả lời:} \dotfill
\end{ex}
\begin{ex}
Có bao nhiêu hợp chất đơn chức ứng với công thức phân tử $\ce{C3H6O2}$ đều tác dụng được với dung dịch NaOH?
\hfill \textbf{Trả lời:} \dotfill
\end{ex}
\vspace{4pt}
\noindent
{\color{mainGreen}\textbf{\Large B. TỰ LUẬN (3,0 điểm).}}\\
\textit{\small Thí sinh trình bày chi tiết lời giải từ câu 1 đến câu 3 vào giấy làm bài thi.}
\vspace{3pt}

\noindent
\textbf{Câu 1. (1,0 điểm)} Mỡ bò có hàm lượng tristearin $80\%$ về khối lượng, các loại chất béo khác có hàm lượng không đáng kể và một số tạp chất khác. Xà phòng hóa hoàn toàn $222,5\text{ gam}$ mỡ bò trong dung dịch KOH, thu được $m\text{ gam}$ potassium stearate.
\begin{enumerate}[label=\alph*),leftmargin=20pt,itemsep=1pt]
\item Viết phương trình hóa học?
\item Tính giá trị của $m$? Cho rằng các tạp chất không tham gia phản ứng với KOH.
\end{enumerate}

\vspace{1pt}
\noindent
\multido{}{7}{\noindent\makebox[\linewidth]{\dotfill}\\[14pt]}

\vspace{2pt}
\noindent
\textbf{Câu 2. (1,0 điểm)} Hoàn thành các phương trình phản ứng sau:
\begin{enumerate}[label=\alph*),leftmargin=20pt,itemsep=1pt]
\item Glucose + dung dịch $\ce{Br2}$
\item Glycine + NaOH
\item Propan-1-amine + HCl
\item Methyl acetate + KOH ($t^\circ$)
\end{enumerate}

\vspace{1pt}
\noindent
\multido{}{7}{\noindent\makebox[\linewidth]{\dotfill}\\[14pt]}

\vspace{2pt}
\noindent
\textbf{Câu 3. (1,0 điểm)} Cho các pin điện hoá và sức điện động chuẩn tương ứng:
\begin{center}
\renewcommand{\arraystretch}{1.2}
\begin{tabular}{|c|c|c|c|}
\hline
\rowcolor{tableGreen}
\textbf{Pin điện hóa} & \textbf{Ni -- Cu} & \textbf{Zn -- Cu} & \textbf{Ni -- Sn} \\
\hline
\textbf{Sức điện động chuẩn (V)} & 0,59 & 1,10 & 0,11 \\
\hline
\end{tabular}
\end{center}
\begin{enumerate}[label=\alph*),leftmargin=20pt,itemsep=1pt]
\item Viết phương trình ion diễn ra khi pin Zn -- Sn hoạt động.
\item Sức điện động chuẩn của pin Zn -- Sn bằng bao nhiêu V(volt)?
\end{enumerate}

\vspace{1pt}
\noindent
\multido{}{8}{\noindent\makebox[\linewidth]{\dotfill}\\[14pt]}

\vfill
\begin{center}
\textbf{---------- HẾT ----------}
\end{center}
\end{document}
